% =============================================================================
%                    TEMPLATE VOOR CURSUSSAMENVATTING
%                    KU Leuven - School Macros v4.0
% =============================================================================
% Dit is een starttemplate voor nieuwe documenten.
% Kopieer dit bestand en pas het aan voor je eigen cursus.
%
% Tip: Compileer TWEEMAAL voor correcte formularium en symbolenlijst!
% =============================================================================

\documentclass[a4paper,11pt]{article}
\usepackage{school-macros}
\enablecompactmode

% v4.0 Optioneel: Compact Mode voor cheatsheets
% \enablecompactmode

% =============================================================================
%                           DOCUMENT INFO
% =============================================================================
% Pas deze gegevens aan voor je cursus:
\newcommand{\vaknaam}{Object gericht-programmeren}
\newcommand{\auteur}{Ruben Ryckaert}

\title{\vaknaam{} --- Synthax Java}
\author{\auteur}

\date{\today}


% =============================================================================
%                              DOCUMENT
% =============================================================================
\begin{document}
\maketitle

 
\begin{abstract}
Deze samenvattingen zijn beschikbaar op GitHub waar je kunt bijdragen aan verbeteringen: \url{https://github.com/Eggmansmile/Samenvattingen-Ku-leuven-Industriele-ingenieurs}

Je bent helemaal vrij en ik raad je zelf aan om info toe te voegen aan de documenten als je de uitleg niet goed vindt.
Je leert je vak en je raakt dan ook bekend met LaTeX wat handig is voor later.
\end{abstract}


\section{import statements}
Om klassen van andere packages te gebruiken, moeten deze eerst geïmporteerd worden met een import statement.\\

\begin{codeblock}[Import utils]{java}
    import java.util.*; // importeert alle klassen van de package java.util
    import java.util.ArrayList; // importeert enkel de ArrayList klasse
\end{codeblock}

\section{typische Java klasse structuur}
Een typische Java klasse heeft de volgende structuur:
\begin{codeblock}[Hoe java doc starten?]{java}
    package packageNaam; // optioneel

    import package1.Klasse1; // optioneel
    import package2.*; // optioneel

    public class KlasseNaam { // klasse definitie

        // attributen
        private int attribuut1;
        private String attribuut2;

        // constructor
        public KlasseNaam(int attribuut1, String attribuut2) {
            this.attribuut1 = attribuut1;
            this.attribuut2 = attribuut2;
        }

        // methoden
        public void methode1() {
            // methode implementatie
        }

        public int methode2(String parameter) {
            // methode implementatie
            return 0;
        }
    }
\end{codeblock}



\end{document}
